% ===== PRÉAMBULE CANONIQUE — à coller VERBATIM en tête de CHAQUE .tex =====
\documentclass[11pt,a4paper]{article}
\usepackage[utf8]{inputenc}\usepackage[T1]{fontenc}\usepackage{lmodern}
\usepackage[french]{babel}
\usepackage{amsmath,amssymb,mathtools}\usepackage{icomma}
\usepackage[version=4]{mhchem}          % \ce{...} : équations & formules
\usepackage{siunitx}\sisetup{locale=FR} % \qty{}{}, \si{}, \ang{}
\usepackage{chemfig}                    % \chemfig{...} : structures développées
\usepackage{xcolor}\usepackage{colortbl}
\usepackage{tikz}\usepackage{pgfplots}\pgfplotsset{compat=1.18}
\usepackage[most]{tcolorbox}\usepackage{enumitem}\usepackage{array,booktabs}
\usepackage[a4paper,margin=2.2cm]{geometry}
\usetikzlibrary{arrows.meta,calc,angles,quotes,positioning,patterns,backgrounds,shapes.geometric}
\definecolor{primary}{RGB}{222,49,99}\definecolor{primarydark}{RGB}{150,22,55}
\definecolor{defcol}{RGB}{0,114,178}\definecolor{methcol}{RGB}{52,92,148}
\definecolor{excol}{RGB}{0,135,90}\definecolor{attncol}{RGB}{200,40,55}
\tcbset{boxbase/.style={enhanced,breakable,boxrule=0.9pt,arc=2.5pt,
  left=3mm,right=3mm,top=2.3mm,bottom=2.3mm,fonttitle=\bfseries,coltitle=white}}
% --- boîtes TUTORIEL (noms EXACTS, ne pas renommer) ---
\newtcolorbox{cle}[1][]{boxbase,colback=defcol!6,colframe=defcol,
  title={\textsc{À retenir}\ifx\relax#1\relax\relax\else\textmd{ : \itshape #1}\fi}}
\newtcolorbox{methode}[1][]{boxbase,colback=methcol!7,colframe=methcol,sharp corners,
  title={\textsc{Méthode}\ifx\relax#1\relax\relax\else\textmd{ --- \itshape #1}\fi}}
\newtcolorbox{exemplebox}[1][]{boxbase,colback=excol!5,colframe=excol,
  title={\textsc{Exemple}\ifx\relax#1\relax\relax\else\textmd{ : \itshape #1}\fi}}
\newtcolorbox{piege}[1][]{boxbase,colback=attncol!6,colframe=attncol,
  title={\textsc{Piège}\ifx\relax#1\relax\relax\else\textmd{ : \itshape #1}\fi}}
\newtcolorbox{remarque}[1][]{boxbase,colback=black!4,colframe=black!45,
  title={\textsc{Remarque}\ifx\relax#1\relax\relax\else\textmd{ : \itshape #1}\fi}}
% --- boîtes EXERCICE / CORRIGÉ (noms EXACTS, auto-comptées) ---
\newtcolorbox[auto counter]{exo}[1][]{boxbase,colback=primary!4,colframe=primary,
  title={\textsc{Exercice}~\thetcbcounter\ifx\relax#1\relax\relax\else\textmd{ --- \itshape #1}\fi}}
\newtcolorbox[auto counter]{corrige}[1][]{boxbase,colback=excol!4,colframe=excol,
  title={\textsc{Corrigé}~\thetcbcounter\ifx\relax#1\relax\relax\else\textmd{ --- \itshape #1}\fi}}
\newcommand{\R}{\mathbb{R}}\newcommand{\N}{\mathbb{N}}\newcommand{\Z}{\mathbb{Z}}\newcommand{\ds}{\displaystyle}
% (un \usepackage{fancyhdr}+en-tête est possible ; il est ignoré côté web)
% =========================================================================
\begin{document}
\begin{exo}[configuration : ouvrir la couche N]
Écris la configuration électronique de chaque atome neutre. \emph{Attention :} la couche externe ne
dépasse pas $8$, ce qui oblige parfois à ouvrir une couche N.
\begin{enumerate}[label=\alph*)]
  \item argon $\ce{Ar}$ ($Z = 18$)
  \item potassium $\ce{K}$ ($Z = 19$)
  \item calcium $\ce{Ca}$ ($Z = 20$)
\end{enumerate}
\end{exo}

\begin{exo}[configuration électronique d'un ion]
Écris la configuration électronique de chaque ion (on part de l'atome neutre, puis on retire ou ajoute
des électrons).
\begin{enumerate}[label=\alph*)]
  \item $\ce{Na^{+}}$ ($Z = 11$)
  \item $\ce{Mg^{2+}}$ ($Z = 12$)
  \item $\ce{O^{2-}}$ ($Z = 8$)
  \item $\ce{Cl^{-}}$ ($Z = 17$)
\end{enumerate}
\end{exo}

\begin{exo}[répartition par couche]
Donne la configuration de l'ion, puis le nombre d'électrons présents sur \emph{chaque} couche et la
couche de valence.
\begin{enumerate}[label=\alph*)]
  \item $\ce{Cl^{-}}$ ($Z = 17$)
  \item $\ce{Ca^{2+}}$ ($Z = 20$)
\end{enumerate}
\end{exo}

\begin{exo}[vrai ou faux : configurations]
Indique si chaque proposition est \textbf{vraie} ou \textbf{fausse}, en corrigeant si besoin.
\begin{enumerate}[label=\alph*)]
  \item La configuration du calcium ($Z=20$) est $(\mathrm{K})^2(\mathrm{L})^8(\mathrm{M})^{10}$.
  \item La couche de valence du phosphore ($Z=15$) est $(\mathrm{L})^5$.
  \item L'ion $\ce{Na^{+}}$ a la même configuration électronique que l'atome de néon ($Z=10$).
  \item Un atome de soufre ($Z=16$) possède $6$ électrons de valence.
\end{enumerate}
\end{exo}

\begin{exo}[symbole de Lewis et célibataires]
Pour chaque atome, donne le nombre d'électrons de valence, le nombre de doublets et de célibataires.
\begin{enumerate}[label=\alph*)]
  \item bore $\ce{B}$ ($Z = 5$)
  \item oxygène $\ce{O}$ ($Z = 8$)
  \item silicium $\ce{Si}$ ($Z = 14$)
\end{enumerate}
\end{exo}

\end{document}
